% Source: https://artofproblemsolving.com/community/c3h2516435_why_the_aops_way_is_sometimes_not_intuitive_or_out_of_the_box
% Stable identity: aops-topic-2516435-q01
\textbf{(a)} Since $r$ and $s$ are the roots, $ax^2+bx+c=a(x-r)(x-s)$. Let $m=\dfrac{r+s}{2}$ and $d=\dfrac{s-r}{2}$, so $r=m-d$ and $s=m+d$. Then
\[a(x-r)(x-s)=a\bigl((x-m)+d\bigr)\bigl((x-m)-d\bigr)=a(x-m)^2-ad^2.\]
Since $(x-m)^2\ge0$ with equality only at $x=m$, the function has its extreme value (minimum if $a>0$, maximum if $a<0$) exactly at $x=m$. So the vertex is at $x=m=\dfrac{r+s}{2}$, the average of the $x$-intercepts, and the parabola is symmetric about the line $x=m$.

\textbf{(b)} $2x^2-8x-10=2\left(x^2-4x-5\right)=2(x-5)(x+1)$, so $r=-1$, $s=5$ and the vertex has $x=2$; then $y=8-16-10=-18$. The vertex is $\boxed{(2,-18)}$.

\begin{asy}
usepackage("amsmath");
size(5.5cm);
defaultpen(fontsize(10pt));
import graph;
real f(real x){return 2x^2-8x-10;}
pair T(real x,real y){return (x,0.12*y);}
draw(T(-2.3,0)--T(6.4,0),Arrow(5)); draw(T(0,-21)--T(0,8),Arrow(5));
label("$x$",T(6.4,0),E); label("$y$",T(0,8),N);
guide g; for(int i=0;i<=100;++i){real x=-1.6+7.2*i/100; g=g--T(x,f(x));}
draw(g,blue+linewidth(0.9));
dot(T(-1,0)); dot(T(5,0)); dot(T(2,-18),red);
label("$(-1,0)$",T(-1,0),SW,fontsize(8pt)); label("$(5,0)$",T(5,0),SE,fontsize(8pt));
label("$(2,-18)$",T(2,-18),S,red+fontsize(8pt));
draw(T(2,-18)--T(2,6),dashed+red); label("$x=2$",T(2,6),N,red+fontsize(8pt));
\end{asy}

\textbf{(c)} The discriminant of $x^2+2x+5$ is $4-20<0$, so the parabola has no $x$-intercepts and there is nothing to average. Completing the square, $x^2+2x+5=(x+1)^2+4$, so the vertex is $\boxed{(-1,4)}$. (The idea of (a) still works with any horizontal line that meets the parabola twice: $y=5$ meets it where $x^2+2x=0$, at $x=0$ and $x=-2$, whose average is $-1$.)
